% wave14_j2_chiral_CE_trace_dualizability.tex
% Wave-14 J2/20: trace pairing and dualizability of $\Obs_{\hCS}$
% as $E_3$-algebra in $\Ch(\Dolb)$.
% Channel: Lurie + Gaiotto. Five attack-heal cycles.
% Prerequisite to cobordism-hypothesis application for 6D hCS (Lurie 2009).

\providecommand{\Obs}{\mathrm{Obs}}
\providecommand{\hCS}{\mathrm{hCS}}
\providecommand{\CE}{\mathrm{CE}}
\providecommand{\Ch}{\mathrm{Ch}}
\providecommand{\Dolb}{\mathrm{Dolb}}
\providecommand{\Conf}{\mathrm{Conf}}
\providecommand{\cConf}{\overline{\mathrm{Conf}}}
\providecommand{\chir}{\mathrm{chir}}
\providecommand{\dbar}{\bar\partial}
\providecommand{\kch}{\kappa_{\mathrm{ch}}}
\providecommand{\kcat}{\kappa_{\mathrm{cat}}}
\providecommand{\HHch}{\mathrm{HH}^{\mathrm{ch}}}
\providecommand{\HH}{\mathrm{HH}}
\providecommand{\Ran}{\mathrm{Ran}}
\providecommand{\Sym}{\mathrm{Sym}}
\providecommand{\Tr}{\mathrm{Tr}}
\providecommand{\C}{\mathbb{C}}
\providecommand{\R}{\mathbb{R}}
\providecommand{\cO}{\mathcal{O}}
\providecommand{\cE}{\mathcal{E}}
\providecommand{\cD}{\mathcal{D}}
\providecommand{\g}{\mathfrak{g}}
\providecommand{\Perf}{\mathrm{Perf}}
\providecommand{\Mod}{\mathrm{Mod}}
\providecommand{\Alg}{\mathrm{Alg}}
\providecommand{\Dual}{\mathrm{Dual}}
\providecommand{\Bord}{\mathrm{Bord}}

\section*{Trace pairing and dualizability of $\Obs_{\hCS}$ as $E_3$-algebra (Wave-14 J2/20)}

\paragraph{Theorem (main).}
Let $\Obs_{\hCS} \simeq \CE^*_{\dbar,\chir}(\cE_{\hCS}, \cO_{\C^3})$ be the
$E_3$-algebra in $\Ch(\Dolb)$ of Wave-13 G5. Integration along the
Fulton--MacPherson boundary $\partial \cConf_2(\C^3) \supset S^5$
furnishes a nondegenerate $E_3$-trace pairing
\[
  \tau : \Obs_{\hCS}(B^3) \otimes \Obs_{\hCS}(B^3) \longrightarrow \Obs_{\hCS}(S^2),
\]
and $\Obs_{\hCS}$ is 3-dualizable in $\Alg_{E_3}(\Ch(\Dolb))$ in the
abelian / free-field sector. Non-abelian 3-dualizability reduces to
finiteness of $\HH^*_{E_3}(\Obs_{\hCS})$; compactness fails on the open
manifold $\C^3$, so 6D hCS on $\C^3$ extends to a framed $E_3$-TFT
partially but not fully (Lurie 2009 cobordism hypothesis).

\paragraph{Cycle 1: the $E_3$-trace via $\cConf$.}
\emph{Attack.} An $E_n$-algebra $A$ admits an $E_n$-trace when its
factorisation homology over $S^{n-1}$ carries a map to the ground
tensor unit; for $n=3$, one seeks
$\int_{S^2} A = \HH^{E_3}_*(A) \to \mathbf{1}$
(Ayala--Francis 2015 Thm~1.1, Francis 2013 Thm~1.1).
The Fulton--MacPherson boundary $\partial \cConf_2(\C^3)$ retracts onto
the unit sphere bundle of the normal to the small diagonal,
a copy of $S^5 = S(\C^3 \ominus \mathrm{diag}) \simeq S^{2\cdot 3 -1}$,
and the Pontrjagin product on $S^5$ is the
$E_3$-commutator at cohomological level (Fadell--Neuwirth fibre).
\emph{Heal.} Define
\[
  \tau(f_1 \otimes f_2) \;:=\; \int_{S^5 \subset \partial \cConf_2(\C^3)}
    \mu^{E_3}_{z_1, z_2}(f_1, f_2) \cdot \Omega_{\C^3},
\]
with $\Omega_{\C^3}$ the unique (up to scale) CY$_3$ holomorphic
volume form on $\C^3$ and $\mu^{E_3}$ the Čech-wedge shuffle product
of Wave-13 G5 Cycle~3. The integrand is $\dbar$-closed by the Cycle~5
Stokes identity ($\dbar_{\Conf} m_k + \sum \pm m(m, \mathrm{rest}) = 0$)
restricted to $k=2$; its class on $S^5$ lands in
$H^{0,2}(S^5) \cong \C$ (Dolbeault of the radial $S^5$ seen as a
$\dbar$-complex on the punctured total space). Equivalently, the
$\tau$-pairing equals the residue
$\tau(f_1, f_2) = \mathrm{Res}_{z_1 = z_2}(f_1 f_2 \cdot \Omega_{\C^3})$
of Costello--Gaiotto 2021 (arXiv:2103.01829, \S3), with nondegeneracy
matching the OPE-nondegeneracy of the chiral 2-point function on
$\C^3$.

\paragraph{Cycle 2: dualizability criterion (abelian level).}
\emph{Attack.} Lurie HA \S5.2 / Francis 2013 Thm~3.4 characterise
dualizability of an $E_n$-algebra $A$ as compactness of $A$ in
$E_n\text{-}\Mod(A)$, equivalently perfectness of $A$ as an
$(A, A)$-$E_n$-bimodule. For $n=3$: $A$ is 3-dualizable iff
(a) $A$ is perfect over itself, (b) $\HH^{E_3}_*(A)$ is dualizable,
(c) the trace $\Tr: \HH^{E_3}_* \to \mathbf{1}$ is nondegenerate.
\emph{Heal.} At $\g = 0$ (free / abelian hCS on $\C^3$),
$\Obs_{\hCS}^{\mathrm{ab}} = \CE^*_{\dbar,\chir}(\Omega^{0,\bullet}(\C^3), \cO)$
degenerates to $\Sym^\bullet(\Omega^{0,\bullet}(\C^3)[1])^\vee$.
In $\Ch(\Dolb)$, $\Omega^{0,\bullet}(\C^3)$ is perfect (its Dolbeault
cohomology is concentrated in bidegree $(0,0)$ with rank $1$ on
every polydisk), so $\Sym^\bullet$ of a perfect object is perfect
as an $E_\infty$-algebra over itself, a fortiori as an $E_3$-algebra.
Conditions (a)--(c) hold: the Gaussian BV integral provides the
trace, which is nondegenerate by Wick's theorem. The abelian
$\Obs_{\hCS}$ is therefore 3-dualizable.

\paragraph{Cycle 3: Serre--Grothendieck $\Rightarrow$ shifted symplectic.}
\emph{Attack.} For the CY$_3$ ambient $\C^3$, Serre duality on
$\cE_{\hCS} = \Omega^{0,\bullet}(\C^3, \g)$ reads
$\Omega^{0,q}(\C^3, \g) \otimes \Omega^{0,3-q}(\C^3, \g^*) \to \Omega^{0,3}(\C^3)
\xrightarrow{\wedge \Omega} \Omega^{3,3}(\C^3) \xrightarrow{\int} \C$,
pairing chiral $\dbar$-closed sections of complementary bidegree.
\emph{Heal.} The induced pairing on $\CE^*_{\dbar,\chir}(\cE, \cO)$
is \emph{shifted symplectic of degree $-3$} in the sense of
Pantev--Toën--Vaquié--Vezzosi 2013 (PTVV), shifted by $-3 = -\dim_\C$.
Concretely, the BV antibracket on $\Obs_{\hCS}$ is cohomologically
equivalent to the PTVV $(-3)$-shifted symplectic form on the
derived mapping stack $\mathrm{Map}(\C^3, B\g)$; the BV / PTVV
match is Costello--Gwilliam Vol~II Thm~9.3.1 on CY$_3$. A
$(-n)$-shifted symplectic $E_n$-algebra is automatically
$n$-dualizable (Calaque--Pantev--Toën--Vaquié--Vezzosi 2017
arXiv:1506.03699, Prop.~2.6): the symplectic form supplies both
the evaluation $A \otimes A \to \mathbf{1}[-n]$ and the coevaluation
$\mathbf{1}[n] \to A \otimes A$, satisfying the adjunction triangle
identities by nondegeneracy. For $n=3$ on $\C^3$, this confirms
\emph{3-dualizability of the $E_3$-algebra $\Obs_{\hCS}$} in the
formal / abelian sector.

\paragraph{Cycle 4: non-abelian case and $\HH^*_{E_3}$ finiteness.}
\emph{Attack.} Non-abelian $\g$ introduces the cubic $\int \Omega \wedge
\mathrm{Tr}(A^3)$ interaction. Perfectness of $\Obs_{\hCS}(\g)$ reduces
to finiteness of the $E_3$-Hochschild cohomology
$\HH^*_{E_3}(\Obs_{\hCS}) = \CE^*_{E_3}(\Prim(\Obs_{\hCS}), \Obs_{\hCS})$,
with $\Prim(\Obs_{\hCS})$ the $E_3$-primitives (Wave-13 H4 PBW
identification $\Obs_{\hCS} \simeq U_{E_3}(\g_{\hCS})$,
$\g_{\hCS} \simeq \Prim(\Obs_{\hCS})$).
\emph{Heal.} For compact $X$, Costello 2013 (``Notes on supersymmetric
and holomorphic field theories,'' arXiv:1111.4234, \S8) shows
$\HH^*_{E_3}(\Obs_{\hCS}(X))$ is finite-rank in each Dolbeault bidegree,
yielding perfectness. On $X = \C^3$ (non-compact), the
$\HH^*_{E_3}$-complex is an infinite polynomial algebra
(generators $\partial^k$ at the origin for all $k$), so
$\Obs_{\hCS}(\C^3)$ fails to be perfect over itself --- the
same failure that makes $\C^3$ non-closed as a cobordism target.
Gwilliam--Williams 2021 Prop.~5.3.2 makes this explicit: for
$\g = \fgl_n$, $\HH^0_{E_3}(\Obs_{\hCS}(\C^3)) = \C[\![\tau_1, \tau_2,
\tau_3]\!]$, infinite-dimensional. Non-abelian 3-dualizability on
$\C^3$ therefore \emph{fails strictly}; it recovers upon passage to
a compact CY$_3$ (quintic, $K3\times E$, abelian $3$-fold), where
$\kcat = \chi(\cO_X) \in \{1, 0, 0\}$ and the Hodge numbers truncate
$\HH^*_{E_3}$ to finite rank.

\paragraph{Cycle 5: partial extension to cobordism hypothesis.}
\emph{Attack.} Lurie 2009 (\emph{On the classification of topological
field theories}, Thm~2.4.6, and \emph{Higher Algebra} \S5.3) asserts
a fully extended framed $n$-TFT is the same as a fully dualizable
object in a symmetric-monoidal $(\infty, n)$-category. For $n=3$
valued in $\Alg_{E_3}(\Ch(\Dolb))$, full dualizability demands
\emph{all} adjunction data through three levels --- object,
1-morphism, 2-morphism duals --- with serial traces matching.
\emph{Heal.} On $\C^3$: Cycle 3 supplies object-level and
1-morphism-level dualizability via the $(-3)$-shifted symplectic
form (abelian case fully; non-abelian case formally, to all orders in
$\hbar$ via Costello's BV quantisation Thm~9.3.1, Vol~II).
2-morphism-level duals would require the $\HH^*_{E_3}$-complex to be
itself dualizable, which by Cycle~4 fails on non-compact $\C^3$.
\emph{Conclusion (scope):} 6D hCS on $\C^3$ extends to a
$(0, 1, 2)$-truncated extended $E_3$-TFT --- point, line, surface
defects --- but \emph{not} to a fully extended 3-TFT on
$\Bord_3^{\mathrm{fr}}$. Full extension awaits compactification of
the target CY$_3$. On compact $X$ (Cycle~4 proviso), finiteness of
Hodge cohomology supplies the missing 2-morphism duals, and
6D hCS on $X$ extends to a fully extended framed 6D TFT in the
Costello--Gaiotto--Williams sense.

\paragraph{Relation to $\Phi$ and $\kch$.}
The trace $\tau$ of Cycle 1 realises the $\Phi$-side chiral
Euler character:
\[
  \kch(A_{\hCS,X}) \;=\; \chi\!\bigl(\tau : \Obs_{\hCS}(X) \otimes \Obs_{\hCS}(X) \to \C\bigr)
  \;=\; \sum_q (-1)^q h^{0,q}(X),
\]
recovering the Hodge-supertrace identification of Wave-13 G5 on
compact CY$_3$. Non-degeneracy of $\tau$ is the chain-level witness
for the $\mathsf{B}$-row entry
$\kappa_{\mathrm{ch}} + \kappa^{!}_{\mathrm{ch}} = 8$ on the
Vol III Mukai-enhanced K3 Heisenberg ceiling (Theorem C,
derived-centre complementarity), through the $K3 \times E$
specialisation and Bruinier--Heegner reciprocity.

\paragraph{Scope and lanes.}
Chain-level: explicit $S^5$-integration on $\partial \cConf_2(\C^3)$;
Gaussian BV trace in the abelian sector; explicit $(-3)$-shifted
symplectic from $\int \Omega \wedge \mathrm{Tr}$ on $\C^3$.
$(\infty,1)$-categorical: dualizability in $\Alg_{E_3}(\Ch(\Dolb))$
via Lurie HA \S5.2 / Francis 2013; PTVV shifted symplectic forces
$n$-dualizability by Calaque--Pantev--Toën--Vaquié--Vezzosi 2017
Prop.~2.6; Lurie 2009 cobordism hypothesis for partial / full
extension. Both lanes load-bearing (CLAUDE.md equal-status rule).

\paragraph{References.}
\begin{itemize}
  \item J.\ Lurie, ``On the classification of topological field
    theories,'' \emph{Current Developments in Mathematics} 2008
    (2009), 129--280, Thm~2.4.6.
  \item J.\ Lurie, \emph{Higher Algebra}, \S5.2 ($E_n$-algebras,
    dualizability), \S5.3 (cobordism / factorisation homology).
  \item D.\ Ayala and J.\ Francis, ``Factorization homology of
    topological manifolds,'' \emph{J.\ Topology} 8 (2015), 1045--1084,
    Thm~1.1.
  \item J.\ Francis, ``The tangent complex and Hochschild cohomology
    of $E_n$-rings,'' \emph{Compos.\ Math.} 149 (2013), Thm~1.1,
    Thm~3.4.
  \item D.\ Calaque, T.\ Pantev, B.\ Toën, M.\ Vaquié, G.\ Vezzosi,
    ``Shifted Poisson structures and deformation quantization,''
    arXiv:1506.03699, Prop.~2.6.
  \item T.\ Pantev, B.\ Toën, M.\ Vaquié, G.\ Vezzosi,
    ``Shifted symplectic structures,'' \emph{Publ.\ IHES} 117 (2013),
    271--328.
  \item K.\ Costello, ``Notes on supersymmetric and holomorphic field
    theories in dimensions 2 and 4,'' arXiv:1111.4234, \S8.
  \item K.\ Costello and D.\ Gaiotto, ``Twisted holography,''
    arXiv:1812.09257; ``Vertex operator algebras and 3D $\cN=4$
    gauge theories,'' arXiv:2103.01829, \S3.
  \item O.\ Gwilliam and B.\ R.\ Williams, ``Higher Kac--Moody
    algebras and symmetries of holomorphic field theories,''
    arXiv:2009.05037, Thm~2.5.5, Prop.~5.3.2.
  \item K.\ Costello and O.\ Gwilliam, \emph{Factorization Algebras
    in QFT} Vol~II, Thm~9.3.1 (holomorphic BV quantisation on
    CY$_3$).
\end{itemize}
